\scanpage{255}
\mspage{205}

déduit de l'inclusion
\[
(\underline{\mathrm{Fib}}_W S,W_S^{\unsure{t|f}})\lhook\joinrel\longrightarrow(\mathrm{Cat}_{/S},W_S).
\]
Le fait que \unsure{(xx)|ceci} soit une équivalence, se prouve justement par le fait de l'existence d'un foncteur $C:\mathrm{Cat}_{/S}\to\mathrm{Cat}_{/S}$ et d'un $i:\mathrm{id}_{\mathrm{Cat}_{/S}}\to C$, ayant les propriétés que \unsure{l'on|on} sait, i.e. du fait que $\unsure{\mathcal{C}|C}\ne\varnothing$. De plus, la donnée d'un objet $C$ de $\mathcal{C}$ donne un moyen commode d'expliciter le quasi-inverse du dit foncteur. Mais ce quasi-inverse \unsure{à|a} une existence canonique \unsure{etc|et} indépendante du choix d'un $C$, i.e. d'un mode de calcul. Le fait que le foncteur composé \inserted{$\unsure{\varphi|\phi}$} dans \unsure{(x)|(*)} puisse se ``relever'' \inserted{(à isom près)} \unsure{en|via} un foncteur
\[
C:\mathrm{Cat}_{/S}\longrightarrow\underline{\mathrm{Fib}}_W S,
\]
et que l'isomorphisme \inserted{entre $\varphi(X)\unsure{\flat|b|!}$ et $\mathrm{hot}_{W,S}(C(X))$} \margin{$\flat=\mathrm{hot}_{W,S}(X)$} \unsure{puisse s'expliciter|en question s'explicite}, lui \unsure{même|aussi}, par une flèche $X$-fonctorielle
\[
i_X:X\longrightarrow C(X)\qquad i_X\in W
\]
dans $\mathrm{Cat}_{/S}$ lui-\unsure{$=$|même} (et non \inserted{seulement} dans le \unsure{localisé} $(\mathrm{Cat}_{/S},W_S^{-1})$,
% archive sheet number: 255
